\scanpage{63}
% archive sheet number: 65

a) Cat. des surfaces conformes à bord $X$, avec comp. connexe \inserted{compacte} choisie $C$ du bord, et une structure toroïdale dessus rectifiable (?) pour sa structure analytique réelle, et

b) Cat. des surfaces conformes à bord \inserted{$X'$}, avec \deleted{\unsure{$D$|$H$|un}} disque pointé $(D,a)$, $D\subset\mathrm{Int}(X')$ marqué, tel que $\partial D$ soit rectifiable dans $X'$ (i.e l'application $\partial D\unsure{-|\to}X'$ rectifiable, pour la structure analytique naturelle de $\partial D$ comme bord de $D$).

De plus, pour $(X',a)$ fixé, chaque disque $D$ autour de $x$ définit un ``disque \deleted{canonique} \inserted{euclidien}'' dans $T_{X',a}$ (au sens de la structure \inserted{vectorielle\unsure{-}} conforme \inserted{de $T_{X',\unsure{\cdot|a}}$}). Comme les disques forment un torseur sous $\mathbb R^{*+}$ (par homothéties), donc un espace de dim 1, il est clair que l'application
